\documentclass[a4paper,10pt]{article}

%Packages
% en haut du main.tex
\usepackage{shellesc}
\write18{mkdir -p build}
\pdfoutput=1

\usepackage{txfonts}
\usepackage{titlesec}
\usepackage[utf8]{inputenc}
\usepackage[T1]{fontenc}
\usepackage[french]{babel}
\usepackage{geometry}
\usepackage{graphicx}
\usepackage[intoc]{nomencl}
\usepackage{multicol}
\usepackage{fancyhdr}
\usepackage{hyperref}
\usepackage{float}
\usepackage{setspace}
\usepackage{caption}
\usepackage{booktabs}
\usepackage{amsmath} % pour les équations

% MISE EN PAGE
\geometry{a4paper, top=18mm, bottom=18mm, textwidth=185mm}
\setlength{\columnsep}{6mm} % espace entre les colonnes
\setlength{\parindent}{0pt} % espace indentation

\titleformat{\section}
{\normalfont\bfseries} % <== style du titre
	{\thesection.}{0.5em}{} % numéro + espacement
\titleformat{\subsection}
	{\normalfont\itshape} % <== style du sous-titre
	{\thesubsection.}{0.5em}{} % numéro + espacement

\pagestyle{fancy}
\renewcommand{\footrulewidth}{0pt} %retire la ligne du bas
\fancyhf{}
\rhead{\includegraphics[scale=0.4]{Images/hepia_simple.png}} %ajoute le logo hepia
\rfoot{\textit{\today}} % date en bas à droite
\cfoot{\thepage} %numéro de page

\captionsetup{
  font=footnotesize,   % taille du texte de la légende
  labelsep=colon       % met un ":" après le numéro
}

% NOMENCLATURE
\makenomenclature
\nomenclature{AFM}{Microscope à Force Atomique}
\nomenclature{$R_a$}{Rugosité moyenne arithmétique}
\nomenclature{$R_q$}{Rugosité RMS (écart-type)}
\nomenclature{$R_{pv}$}{Peak-to-valley (écart max--min)}
\nomenclature{PSPD}{Photodiode sensible à la position (4 quadrants)}
\nomenclature{LFM}{Microscopie en force latérale (friction)}
%\nomenclature{$$}{}

\begin{document} %début du document
 % l'accolade sert à laisser la mise en page en une seule colonne
% TITRES, AUTEURS, ORGANISATION
%
\begin{center}
\vspace*{0mm}
{\large
Rapport de laboratoire\\
\vspace{3mm}
Microscopie à Force Atomique (AFM)
}\\
\vspace{8mm}
% ====================================================================
% TODO : Remplacer par tes noms et emails
% ====================================================================
Guilhem Rozier Vilardell\footnote{\textit{\href{mailto:guilhem.roziervi@hes-so.ch}{MT2, guilhem.roziervi@hes-so.ch}}},
Noah Michelot\footnote{\textit{\href{mailto:noah.michelot@hes-so.ch}{MT2, noah.michelot@hes-so.ch}}},
Mattéo Pozzo Di Borgo\footnote{\textit{\href{mailto:matteo.pozzodi@hes-so.ch}{MT2, matteo.pozzodi@hes-so.ch}}}

%
\vspace{3mm}
\footnotesize{
\textit{Haute école de paysagerie, d'ingénierie et d'architecture, Genève, Suisse}
}\\
\vspace{9mm}
\end{center}
%
\rule{\linewidth}{0.4pt} % ligne horizontale
\onehalfspacing
\textbf{Abstract}\\
% ====================================================================
% TODO : Adapter l'abstract à tes propres résultats (valeurs numériques,
% échantillons étudiés, modes utilisés). L'abstract ci-dessous est un
% template basé sur le rapport de Proietto/Saura.
% ====================================================================
Ce rapport présente une étude de caractérisation de surfaces par microscope
à force atomique (AFM) en mode contact. Trois échantillons ont été analysés :
un wafer de silicium avec couche de SiO$_2$ (échantillon de calibration TGZ3),
un wafer Si/Al déposé à 150\,°C, et un film mince polymère PS/PMMA (50/50)
déposé par spin-coating. Les cartes topographiques ont été traitées avec le
logiciel XEI afin d'extraire les indicateurs de rugosité ($R_a$, $R_q$,
$R_{pv}$) et les dimensions caractéristiques des structures de surface.
Les résultats montrent des rugosités allant de quelques dizaines de nanomètres
à plusieurs micromètres selon l'échantillon.\\
\noindent\textit{Mots-clés :} AFM, mode contact, rugosité, topographie, traitement d'image\\
\rule{\linewidth}{0.4pt} % ligne horizontale

\singlespacing

% TABLE DES MATIERES
\tableofcontents

% NOMENCLATURE
\printnomenclature

\setlength{\parskip}{1pt} % espace entre deux paragraphes

% ====================================================================
% INTRODUCTION
% ====================================================================
\section{Introduction}

L'objectif de ce travail est de caractériser la topographie de surfaces
à l'échelle nanométrique à l'aide d'un microscope à force atomique (AFM).
L'AFM permet d'obtenir une information tridimensionnelle de la
surface avec une résolution latérale de l'ordre de 2\,nm et une résolution
verticale pouvant atteindre 0,1\,nm.

% ====================================================================
% TODO : Adapter les échantillons si les tiens diffèrent de ceux-ci.
% ====================================================================
Trois échantillons ont été étudiés :
\begin{itemize}
    \item un échantillon de calibration TGZ3 (wafer Si + SiO$_2$,
          marches de hauteur nominale connue) ;
    \item un wafer de silicium avec couche d'aluminium déposée à 150\,°C ;
    \item un film mince polymère PS/PMMA (50/50) déposé par spin-coating.
\end{itemize}

Ce rapport décrit le principe de l'AFM en mode contact, la procédure
expérimentale, les résultats obtenus après traitement d'image et une
analyse des sources d'incertitudes.


% ====================================================================
% THÉORIE
% ====================================================================
\section{Théorie}

\subsection{Principe de l'AFM en mode contact}

Le microscope à force atomique utilise un levier (cantilever) souple muni
d'une pointe nanométrique pour sonder la surface d'un échantillon. Le
levier, d'épaisseur typique $0{,}5$--$2\;\mu$m et de longueur $50$--$200\;\mu$m,
est caractérisé par sa constante de raideur $k$ (de l'ordre de $0{,}1$--$10$\,N/m).

Le faisceau d'une diode superluminescente (SLD) est focalisé sur l'extrémité
du levier. Le faisceau réfléchi est dirigé vers une photodiode à 4 quadrants
(PSPD). Les signaux des quadrants permettent de mesurer :
\begin{itemize}
    \item la déflection verticale du levier : signal $A - B = (1+2)-(3+4)$ ;
    \item la torsion latérale (friction, LFM) : signal $C - D = (1+3)-(2+4)$.
\end{itemize}

\subsection{Mode à force constante}

En mode contact à force constante, une boucle de contre-réaction (PID)
ajuste en permanence la position $z$ du scanner piézoélectrique pour
maintenir la déflection du levier --- et donc la force exercée sur
l'échantillon --- constante. Le signal de commande envoyé au scanner $z$
constitue le signal de topographie $z(x,y)$.

% ====================================================================
% TODO : Si tu veux, tu peux ajouter ici un schéma-bloc de
% l'asservissement (figure du protocole, Figure 4). Sinon, supprime
% le \begin{figure} ci-dessous.
% ====================================================================
%\begin{figure}[H]
%\centering
%\includegraphics[width=\linewidth]{Images/schema_AFM.png}
%\caption{Schéma du système d'asservissement de l'AFM en mode contact.}
%\label{fig:schema_afm}
%\end{figure}

\subsection{Paramètres de rugosité}

La rugosité RMS (écart-type des hauteurs) est définie par :
\begin{equation}
    R_q = \sqrt{\frac{1}{N-1}\sum_{i=1}^{N}(z_i - \bar{z})^2}
    \label{eq:Rq}
\end{equation}

La rugosité moyenne arithmétique :
\begin{equation}
    R_a = \frac{1}{N}\sum_{i=1}^{N}|z_i - \bar{z}|
    \label{eq:Ra}
\end{equation}

L'écart peak-to-valley :
\begin{equation}
    R_{pv} = z_{\max} - z_{\min}
    \label{eq:Rpv}
\end{equation}

% ====================================================================
% DESCRIPTION DE L'EXPÉRIENCE
% ====================================================================
\section{Description de l'expérience}

\subsection{Instrument}

L'AFM utilisé est un Park Systems NX7. Le scanner XY a une plage maximale
de $50\;\mu$m ($100\,\%$) et le scanner $z$ une plage de $7{,}5\;\mu$m
($70\,\%$). Le levier utilisé est un Nanoworld CONTR (mode contact).

\subsection{Procédure}

La procédure suivie est celle du protocole \textit{X'Press Manual} :
\begin{enumerate}
    \item Alignement optique : positionnement du spot laser sur l'extrémité
          du levier, centrage sur la PSPD ($A+B > 1$\,V).
    \item Placement de l'échantillon sur le support magnétique.
    \item Pré-positionnement de la tête à $\sim 1$\,mm de la surface.
    \item Configuration : mode Contact, canaux Z Height, Lateral (C-D),
          Error Signal. Flatten = Line, Data Saving = Raw.
    \item Approche automatique.
    \item Vérification du Trace Quality ($> 9{,}0$).
    \item Acquisition de l'image (256$\times$256 pixels par défaut).
    \item Traitement d'image sous XEI : correction de pente ligne par
          ligne (flatten, ordre 1), mesure de profils, extraction des
          rugosités.
\end{enumerate}

% ====================================================================
% RÉSULTATS ET ANALYSE
% ====================================================================
\section{Résultats et analyse}

\subsection{Échantillon 1 : calibration TGZ3}

% ====================================================================
% TODO : Remplacer les images par les tiennes.
% Mettre tes fichiers dans le dossier Images/ et adapter les noms.
% ====================================================================

%\begin{figure}[H]
%\centering
%\includegraphics[width=\linewidth]{Images/ech1_topo.png}
%\caption{Image topographique de l'échantillon de calibration TGZ3
%         ($10 \times 10\;\mu$m$^2$).}
%\label{fig:ech1_topo}
%\end{figure}

L'image topographique de l'échantillon de calibration montre des marches
périodiques gravées dans la couche de SiO$_2$. Une coupe perpendiculaire
aux marches permet d'extraire les grandeurs suivantes :

% ====================================================================
% TODO : Remplacer les valeurs ci-dessous par tes mesures réelles.
% Les valeurs ici proviennent du rapport Proietto/Saura.
% ====================================================================
\begin{itemize}
    \item Hauteur de marche : $531{,}2$\,nm
    \item Période (distance entre deux phases montantes) : $3{,}283\;\mu$m
\end{itemize}

%\begin{figure}[H]
%\centering
%\includegraphics[width=\linewidth]{Images/ech1_profil.png}
%\caption{Profil de hauteur le long de la coupe rouge (échantillon 1).}
%\label{fig:ech1_profil}
%\end{figure}

Les rugosités mesurées sur chaque plan sont résumées dans le
tableau~\ref{tab:ech1_rugo}.

% ====================================================================
% TODO : Remplacer par tes valeurs.
% ====================================================================
\begin{table}[H]
\centering
\caption{Rugosité de l'échantillon de calibration TGZ3.}
\label{tab:ech1_rugo}
\begin{tabular}{lcc}
\toprule
 & Plan inférieur & Plan supérieur \\
\midrule
$R_a$ / nm & 28{,}3 & 33{,}0 \\
$R_q$ / nm & 33{,}5 & 39{,}9 \\
\bottomrule
\end{tabular}
\end{table}

% ====================================================================
% TODO : Commenter la cohérence avec la valeur nominale de la TGZ3
% (hauteur de marche nominale = 518 nm ± 5 nm selon le fabricant MikroMasch).
% Si ta valeur mesurée diffère, discuter l'écart.
% ====================================================================

\subsection{Échantillon 2 : Si/Al déposé à 150\,°C}

% ====================================================================
% TODO : Remplacer les images et les valeurs.
% ATTENTION : dans le rapport Proietto/Saura, les rugosités sont en mV
% et non en nm, ce qui est probablement une erreur de calibration ou
% d'unité dans XEI. Vérifier dans tes mesures que les unités sont
% bien en nm ou µm, PAS en mV.
% ====================================================================

%\begin{figure}[H]
%\centering
%\includegraphics[width=\linewidth]{Images/ech2_topo.png}
%\caption{Image topographique de l'échantillon Si/Al
%         ($5 \times 5\;\mu$m$^2$).}
%\label{fig:ech2_topo}
%\end{figure}

L'image topographique révèle une surface granulaire avec des îlots
d'oxyde en surépaisseur. Les zones claires correspondent aux sommets
des grains, les zones sombres aux creux entre les grains.

%\begin{figure}[H]
%\centering
%\includegraphics[width=\linewidth]{Images/ech2_profil.png}
%\caption{Profil de hauteur le long de la coupe (échantillon 2).}
%\label{fig:ech2_profil}
%\end{figure}

% ====================================================================
% TODO : Remplacer ces valeurs. Vérifier les UNITÉS (nm, pas mV !).
% ====================================================================
Les paramètres de rugosité pour une zone de $5 \times 5\;\mu$m$^2$ :
\begin{itemize}
    \item $R_a = $ \textbf{[À COMPLÉTER]} nm
    \item $R_q = $ \textbf{[À COMPLÉTER]} nm
    \item $R_{pv} = $ \textbf{[À COMPLÉTER]} nm
\end{itemize}

% ====================================================================
% TODO : Si tu as fait des mesures de taille de grain, les ajouter ici.
% ====================================================================


\subsection{Échantillon 3 : film PS/PMMA}

% ====================================================================
% TODO : Adapter le mode de mesure utilisé. Dans le rapport Proietto,
% c'est un mode "force constante" avec indentation pixel par pixel
% (probablement PinPoint ou Force Volume). Vérifier quel mode tu as
% utilisé et adapter le texte.
% ====================================================================

Pour cet échantillon polymère, un mode d'acquisition par indentation
ponctuelle a été utilisé : la pointe s'arrête à chaque pixel, pénètre
la surface, puis avance au point suivant. Cela permet d'obtenir
simultanément la topographie et une carte d'adhésion.

%\begin{figure}[H]
%\centering
%\includegraphics[width=\linewidth]{Images/ech3_topo.png}
%\caption{Image topographique du film PS/PMMA
%         ($30 \times 30\;\mu$m$^2$).}
%\label{fig:ech3_topo}
%\end{figure}

L'image topographique met en évidence une craquelure traversant la zone
observée, ainsi que des amas de polystyrène en surface (zones blanchâtres).

% ====================================================================
% TODO : Remplacer par tes valeurs mesurées.
% ====================================================================

\begin{table}[H]
\centering
\caption{Rugosité du film PS/PMMA ($30 \times 30\;\mu$m$^2$).}
\label{tab:ech3_rugo}
\begin{tabular}{lccc}
\toprule
Zone & $R_a$ / $\mu$m & $R_q$ / $\mu$m & $R_{pv}$ / $\mu$m \\
\midrule
Sans craquelure & 0{,}203 & 0{,}267 & 1{,}729 \\
Avec craquelure & 0{,}563 & 0{,}657 & 2{,}399 \\
\bottomrule
\end{tabular}
\end{table}

La craquelure augmente significativement les paramètres de rugosité :
$R_q$ passe de $0{,}267$ à $0{,}657\;\mu$m (facteur $\times 2{,}5$) et
$R_{pv}$ de $1{,}729$ à $2{,}399\;\mu$m. L'écart de force d'adhésion
entre les domaines PS et PMMA est d'environ 100\,nN, mettant en évidence
le contraste chimique entre les deux polymères.

% ====================================================================
% TODO : Si tu as une image de la carte d'adhésion, l'ajouter ici.
% ====================================================================

%\begin{figure}[H]
%\centering
%\includegraphics[width=\linewidth]{Images/ech3_adhesion.png}
%\caption{Carte d'adhésion du film PS/PMMA.}
%\label{fig:ech3_adhesion}
%\end{figure}


% ====================================================================
% DISCUSSION / SOURCES D'INCERTITUDE
% ====================================================================
\section{Discussion}

% ====================================================================
% TODO : Adapter cette section à tes observations réelles.
% ====================================================================

Plusieurs sources d'incertitude influencent les mesures :
\begin{itemize}
    \item \textbf{Rayon de courbure de la pointe :} une pointe émoussée
          élargit les structures latéralement (convolution pointe--surface)
          et sous-estime la profondeur des creux étroits.
    \item \textbf{Vitesse de balayage et gain :} un gain trop faible
          empêche l'asservissement de suivre fidèlement les structures
          raides ; un gain trop élevé provoque des oscillations du
          scanner $z$.
    \item \textbf{Traitement d'image :} l'ordre du polynôme choisi pour
          le flatten ligne par ligne peut artificiellement supprimer des
          structures réelles (ordre trop élevé) ou laisser des artefacts
          (ordre trop bas).
    \item \textbf{Dérive thermique :} sur des acquisitions longues
          (images haute résolution), la dérive piézoélectrique peut
          affecter la reproductibilité des mesures.
\end{itemize}

% ====================================================================
% TODO : Discuter la comparaison avec la valeur nominale de la TGZ3.
% Par exemple : "La hauteur de marche mesurée (531 nm) diffère de
% la valeur nominale (518 nm) de 2,5 %, ce qui pourrait s'expliquer
% par..."
% ====================================================================


% ====================================================================
% CONCLUSION
% ====================================================================
\section{Conclusion}

% ====================================================================
% TODO : Rédiger ta propre conclusion en résumant tes résultats
% principaux (valeurs numériques) et ce que tu en retires.
% ====================================================================

Ce laboratoire a permis de se familiariser avec l'utilisation de l'AFM
en mode contact et le traitement d'images sous XEI. Les mesures sur
l'échantillon de calibration TGZ3 ont validé le bon étalonnage du
scanner ($h_{\text{marche}} = 531$\,nm). L'échantillon Si/Al présente
une morphologie granulaire avec des îlots d'oxyde en surface. Enfin,
le film PS/PMMA montre un contraste de propriétés mécaniques entre les
deux phases polymères, exploitable via la carte d'adhésion.

L'AFM se confirme comme un outil métrologique puissant, offrant une
information 3D calibrée à l'échelle nanométrique et la capacité de
mesurer simultanément plusieurs propriétés physiques de la surface.


% ====================================================================
% ANNEXE
% ====================================================================
\appendix

\section{Récapitulatif des paramètres d'acquisition}

% ====================================================================
% TODO : Remplir ce tableau avec tes paramètres réels.
% ====================================================================
\begin{table}[H]
\centering
\caption{Paramètres d'acquisition pour chaque échantillon.}
\label{tab:params}
\begin{tabular}{lccc}
\toprule
Paramètre & Éch. 1 & Éch. 2 & Éch. 3 \\
\midrule
Mode & Contact & Contact & Force constante \\
Taille de scan / $\mu$m$^2$ & $10 \times 10$ & $5 \times 5$ & $30 \times 30$ \\
Pixels & $256 \times 256$ & $256 \times 256$ & $256 \times 256$ \\
Scan rate / Hz & 1{,}0 & 1{,}0 & -- \\
Gain & 1 & 1 & -- \\
Flatten (XEI) & Ligne, ordre 1 & Ligne, ordre 1 & Ligne, ordre 1 \\
\bottomrule
\end{tabular}
\end{table}


\section{Figures et Tables}

% ====================================================================
% TODO : Décommenter et ajouter ici les figures supplémentaires
% (vues 3D, profils de coupe, etc.) que tu n'as pas placées dans
% le corps du texte.
% ====================================================================

% Exemple :
%\begin{figure}[H]
%\centering
%\includegraphics[width=0.8\linewidth]{Images/ech1_3D.png}
%\caption{Vue 3D de l'échantillon de calibration TGZ3.}
%\label{fig:ech1_3D}
%\end{figure}

%\begin{figure}[H]
%\centering
%\includegraphics[width=0.8\linewidth]{Images/ech2_3D.png}
%\caption{Vue 3D de l'échantillon Si/Al.}
%\label{fig:ech2_3D}
%\end{figure}

%\begin{figure}[H]
%\centering
%\includegraphics[width=0.8\linewidth]{Images/ech3_3D.png}
%\caption{Vue 3D du film PS/PMMA.}
%\label{fig:ech3_3D}
%\end{figure}

\listoffigures
\listoftables


\end{document}
